% Source: Weinberg GR, physical PDF pp. 163--165
%         (printed pp. 138--140).
% Coverage: Section 6.4, Gravitation versus Curvilinear Coordinates;
%           equations (6.4.1)--(6.4.10).
% Figures/tables/footnotes: none.
% Status: source-reviewed and compile-clean.
% Uncertainties: none.

\subsection{Gravitation versus Curvilinear Coordinates}\label{sec:6.4}

Suppose that we are presented with a metric tensor \(g_{\mu\nu}(x)\) that is
not just a constant. How can we tell if the space is really permeated by a
gravitational field, or if \(g_{\mu\nu}\) merely represents the metric
\(\eta_{\alpha\beta}\) of special relativity written in curvilinear
coordinates? In other words, how can we tell whether there is a set of
Minkowskian coordinates \(\xi^\alpha(x)\) that everywhere satisfy the
conditions
\begin{equation}
  \eta^{\alpha\beta}
  =
  g^{\mu\nu}
  \frac{\partial\xi^\alpha(x)}{\partial x^\mu}
  \frac{\partial\xi^\beta(x)}{\partial x^\nu}?
  \tag{6.4.1}\label{eq:6.4.1}
\end{equation}
Note that the equivalence principle only says that at every point \(X\) we
can find locally inertial coordinates \(\xi_X^\alpha(x)\) that satisfy
Eq.~\eqref{eq:6.4.1} in an infinitesimal neighborhood of \(X\); what we are
asking now is whether we can find one set of coordinates \(\xi^\alpha(x)\)
that satisfy Eq.~\eqref{eq:6.4.1} everywhere. For example, given the metric
coefficients
\begin{equation}
  g_{tt}=-1,
  \qquad
  g_{rr}=1,
  \qquad
  g_{\theta\theta}=r^2,
  \qquad
  g_{\phi\phi}=r^2\sin^2\theta,
  \tag{6.4.2}\label{eq:6.4.2}
\end{equation}
we know that there is a set of \(\xi^\alpha\) satisfying
Eq.~\eqref{eq:6.4.1}, that is,
\begin{equation}
  \xi^0=t,
  \qquad
  \xi^1=r\sin\theta\cos\phi,
  \qquad
  \xi^2=r\sin\theta\sin\phi,
  \qquad
  \xi^3=r\cos\theta,
  \tag{6.4.3}\label{eq:6.4.3}
\end{equation}
but how could we have told that Eq.~\eqref{eq:6.4.2} was really equivalent to
the Minkowski metric \(\eta_{\alpha\beta}\) if we were not clever enough to
have recognized it as simply \(\eta_{\alpha\beta}\) in spherical polar
coordinates? Or, on the other hand, if we change \(g_{rr}\) in
Eq.~\eqref{eq:6.4.2} to an arbitrary function of \(r\), how can we tell that
this really represents a gravitational field; that is, how can we tell that
Eq.~\eqref{eq:6.4.1} now has no solution?

The answer is contained in the following theorem. The necessary and
sufficient conditions for a metric \(g_{\mu\nu}(x)\) to be equivalent to the
Minkowski metric \(\eta_{\alpha\beta}\) [in the sense that there is a
transformation \(x\to\xi\) satisfying Eq.~\eqref{eq:6.4.1}] are, first, that
the curvature tensor calculated from \(g_{\mu\nu}\) must everywhere vanish,
\begin{equation}
  R^\rho{}_{\sigma\mu\nu}=0,
  \tag{6.4.4}\label{eq:6.4.4}
\end{equation}
and, second, that at some point \(X\) the matrix \(g^{\mu\nu}(X)\) has three
positive and one negative eigenvalues.

The necessity of these two conditions is obvious. Suppose that we can find a
coordinate system \(\xi^\alpha(x)\) satisfying Eq.~\eqref{eq:6.4.1}. In this
coordinate system the metric is \(\eta_{\alpha\beta}\), all components of the
affine connection vanish, and hence the Riemann tensor vanishes. But the
vanishing of a tensor is an invariant statement, so the curvature tensor must
have vanished in the original \(x^\mu\) coordinate system. Also, we have
already noted in Section~\ref{sec:3.6} that a ``congruence'' like
Eq.~\eqref{eq:6.4.1} requires \(\eta^{\alpha\beta}\) and \(g^{\mu\nu}\) to
have the same numbers of positive, negative, and zero eigenvalues everywhere.

To prove the sufficiency of Eq.~\eqref{eq:6.4.4} for the existence of a
system of ``everywhere inertial'' coordinates \(\xi^\alpha(x)\) satisfying
Eq.~\eqref{eq:6.4.1}, we shall actually construct the \(\xi^\alpha(x)\).
First we note that at any point \(X\) we may find a matrix
\(d^\alpha{}_\mu\) for which
\begin{equation}
  \eta^{\alpha\beta}
  =
  g^{\mu\nu}(X)d^\alpha{}_\mu d^\beta{}_\nu.
  \tag{6.4.5}\label{eq:6.4.5}
\end{equation}
(For, since \(g^{\mu\nu}(X)\) is a symmetric matrix, we can find an orthogonal
matrix \(O^\alpha{}_\mu\) for which the matrix \(OgO^{\mathsf T}\) is
diagonal, that is, for which
\[
  O^\alpha{}_\mu g^{\mu\nu}O^\beta{}_\nu=D^{\alpha\beta},
  \qquad
  D^{\alpha\beta}
  =
  \begin{cases}
    D^\alpha,&\alpha=\beta,\\
    0,&\alpha\ne\beta.
  \end{cases}
\]
We are assuming that three of the eigenvalues \(D^\alpha\) are positive and
one negative, and can always label the rows of \(O^\alpha{}_\mu\) so that
\(D^1,D^2,D^3\) are positive and \(D^0\) is negative. Rescaling the three
spatial rows by \(1/\sqrt{D^i}\) and the temporal row by
\(1/\sqrt{-D^0}\) then gives a matrix \(d^\alpha{}_\mu\) satisfying
Eq.~\eqref{eq:6.4.5}.)

Next, we define quantities \(D^\alpha{}_\mu(x)\) by the differential equations
\begin{equation}
  \partial_\nu D^\alpha{}_\mu
  =
  \Gamma^\lambda{}_{\mu\nu}D^\alpha{}_\lambda,
  \tag{6.4.6}\label{eq:6.4.6}
\end{equation}
with the initial condition
\begin{equation}
  D^\alpha{}_\mu=d^\alpha{}_\mu
  \qquad\text{at }x=X.
  \tag{6.4.7}\label{eq:6.4.7}
\end{equation}
We showed in the last section that such equations can always be solved,
providing that the curvature tensor vanishes. (The quantities
\(D^\alpha{}_\mu\) are to be thought of as four covariant vectors
\(D^0{}_\mu,D^1{}_\mu,D^2{}_\mu,D^3{}_\mu\), rather than as a single
tensor.) Since \(\partial_\nu D^\alpha{}_\mu\) is symmetric in \(\mu\) and
\(\nu\), we can write the vectors \(D^\alpha{}_\mu\) as gradients of scalars,
which we define to be the locally inertial coordinates \(\xi^\alpha(x)\):
\begin{equation}
  \partial_\mu\xi^\alpha
  =
  D^\alpha{}_\mu,
  \tag{6.4.8}\label{eq:6.4.8}
\end{equation}
with initial values \(\xi^\alpha(X)\) some arbitrary constants. To see that
these \(\xi\) coordinates do satisfy Eq.~\eqref{eq:6.4.1}, note first that
\begin{equation}
  \partial_\rho
  \left(
    g^{\mu\nu}D^\alpha{}_\mu D^\beta{}_\nu
  \right)
  =0.
  \tag{6.4.9}\label{eq:6.4.9}
\end{equation}
This can be verified by direct calculation or, more simply, by noting that
Eq.~\eqref{eq:6.4.6} just says that
\(\nabla_\rho D^\alpha{}_\mu\) vanishes and
\(\nabla_\rho g^{\mu\nu}\) also vanishes, so the covariant derivative of
\(g^{\mu\nu}D^\alpha{}_\mu D^\beta{}_\nu\) vanishes. But since this quantity
is a scalar, this means that its ordinary derivative vanishes. Equations
\eqref{eq:6.4.7} and \eqref{eq:6.4.5} show that
\(g^{\mu\nu}D^\alpha{}_\mu D^\beta{}_\nu\) equals
\(\eta^{\alpha\beta}\) at \(x=X\), so since it is a constant this holds
everywhere:
\begin{equation}
  \eta^{\alpha\beta}
  =
  g^{\mu\nu}D^\alpha{}_\mu D^\beta{}_\nu
  \qquad\text{(all \(x\)).}
  \tag{6.4.10}\label{eq:6.4.10}
\end{equation}
Equation~\eqref{eq:6.4.1} follows immediately from
Eqs.~\eqref{eq:6.4.8} and \eqref{eq:6.4.10}.
